\chapter{Pendahuluan}\label{cha:intro}

Di era digitalisasi pendidikan modern (berdasarkan standar ISO 9241-210 \cite{iso2019ergonomics} mengenai interaksi manusia dan sistem), keberadaan perpustakaan digital (\gls{digilib}) telah menjadi tulang punggung aktivitas akademik mahasiswa. Bagi institusi pendidikan tinggi seperti IIB Darmajaya, platform \gls{digilib} bukan sekadar gudang repositori jurnal atau skripsi, melainkan pusat interaksi literatur ilmiah. Sayangnya, banyak pengembangan platform akademik yang masih terlalu berfokus pada kelengkapan fitur teknis di \textit{back-end}, namun mengabaikan aspek kenyamanan pengguna (\gls{ux}) di \textit{front-end}. Akibatnya, antarmuka yang membingungkan atau alur pencarian referensi yang berbelit-belit seringkali membuat mahasiswa enggan memanfaatkan layanan ini secara maksimal dan lebih memilih mencari referensi di platform eksternal.

Penelitian-penelitian terdahulu mengenai evaluasi sistem informasi perpustakaan seringkali masih bertumpu pada pengujian fungsionalitas dasar (\textit{usability testing}) atau menggunakan kuesioner kepuasan layanan yang konvensional. Metode evaluasi tradisional ini memiliki kelemahan mendasar, yakni hanya mengukur apakah sebuah sistem "berfungsi dan bisa digunakan", tanpa menilai secara mendalam bagaimana perasaan atau pengalaman subjektif pengguna saat berinteraksi dengan sistem tersebut. Sebuah aplikasi yang bebas dari \textit{error} sistem (\textit{bug}) belum tentu memberikan pengalaman yang menyenangkan atau memotivasi penggunanya. Di sisi lain, beberapa metode survei evaluasi \gls{ux} seringkali terlalu panjang dan kompleks, sehingga menyulitkan proses pengumpulan data yang objektif dari responden mahasiswa.

Merespons keterbatasan tersebut, penelitian ini mengusulkan pendekatan evaluasi menyeluruh menggunakan kerangka kerja \gls{ueq} \cite{schrepp2014applying}. Jika evaluasi konvensional hanya melihat aspek pragmatis, pendekatan \gls{ueq} dalam penelitian ini dirancang untuk membedah interaksi pengguna secara komprehensif melalui enam dimensi metrik: \textit{Attractiveness} (daya tarik), \textit{Perspicuity} (kejelasan), \textit{Efficiency} (efisiensi), \textit{Dependability} (ketepatan), \textit{Stimulation} (stimulasi), dan \textit{Novelty} (kebaruan). Pemilihan metode \gls{ueq} didasarkan pada keunggulannya sebagai alat ukur terstandarisasi yang mampu menyeimbangkan aspek kualitas fungsional (pragmatis) dan aspek emosional (hedonis) secara cepat dan reliabel dengan metode analisis kuantitatif.

Kontribusi utama penelitian ini terletak pada dua aspek strategis: (1) Penerapan pengukuran \gls{ux} yang komprehensif untuk mendeteksi secara presisi titik-titik kelemahan desain (\textit{pain points}) pada antarmuka \gls{digilib} IIB Darmajaya tanpa bergantung pada asumsi pengembang semata; dan (2) Penyediaan landasan evaluasi dan rekomendasi perbaikan berbasis data (\textit{data-driven}) yang konkret bagi pengelola sistem informasi kampus, sehingga perombakan sistem di masa depan dapat lebih terarah dalam menciptakan ekosistem akademik digital yang ramah pengguna.